% R. Brown's translations of two letters to L. Breen, with his note and comments (cote 134-2): archive cote 134-2, site page 123 (PDF page 124).
% translation of the letter of 17/2/1975. Draft compared: none.
% AI-transcribed by Claude Opus 5.5 (claude-opus-5-5, adaptive thinking, effort high), Message Batches API, batch brown_full (msgbatch_01Md3i5yK1PxkpJn5GHfEExE), 2026-11-13; not yet checked by eye.
\scanpage{134-2/123}
\mspage{- 2 -}

In the case $n = 1$ , one knows that one can recover (up to equivalence)
the 1-groupoid $C_1$ from the 1-category $\mathit{Hom}(C_1, (\mathrm{Ens}))$ of the functors into
(Ens) = (0 - Cat) (i.e. the ``local systems'' on $C_1$) - which is a topos, called
``multigalois'' - like the category of ``fibred functors'' on the above topos,
i.e. the opposite category to the category of points of this topos (which
is none other than the \textit{classifying topos} of $C_1$) . To make precise for
arbitrary $n$ the way in which the homotopy n-type of a topos $X$ (supposed
for simplicity to be locally homotopically trivial in dim $\leq n$ ) i.e. its
fundamental n-groupoid $C_n$ , can be expressed in terms of the n-category of
``local (n-1)-systems on $X$'' i.e. of the locally constant (n-1)-champs on $X$ ,
and to elucidate completely the hypothetical statement (B) above, it is
necessary to make explicit how an n-groupoid $C_n$ can be recovered, up to
n-equivalence, from the knowledge of the n-category
\[ \mathbb{C}_n = (n\text{-}\mathit{Hom})(C_n , ((n-1)\text{-}\mathrm{Cat})) \]
of local (n-1)-systems on $C_n$ . One would like to say that $C_n$ is the category
of ``fibred n-functors'' on $\mathbb{C}_n$ , i.e. of n-functors $\mathbb{C}_n \rightarrow ((n-1)\text{-}\mathrm{Cat})$ having
certain exactness properties (for $n = 1$ ; this is the condition of being
the inverse image functor for a morphism of toposes, i.e. to commute with
arbitrary $\underrightarrow{\ell im}$ and with finite $\underleftarrow{\ell im}$ ....). It is this which makes real
the fear, expressed in my preceding letter, that one ends by falling upon
the notion of n-topos and of morphisms of these! $\mathbb{C}_n$ will be an n-topos,
(called the ``classifying n-topos'' of the n-groupoid $C_n$) , ((n-1)-Cat) will
be the n-topos of points, and $C_n$ will be interpreted modulo n-equivalence
as the n-category of ``n-points'' of the classifying n-topos $\mathbb{C}_n$ . Brr!

If one hopes to be able to define a good old classifying (1)-topos for
an n-groupoid $C_n$ , as solution of a universal problem, I can see only how
\margin{\unsure{this} seems \unsure{suspect}} to recover the following universal problem: for every topos $T$ , consider
$\mathit{Hom}(\Pi_n(T), C_n)$ . This is an n-category, but take from it the truncated
1-category $\tau_1 \mathit{Hom}(\Pi_n(T), C_n)$ . For variable $T$ , one wants to 2-represent
the contravariant 2-functor $(\mathrm{Top})^{\circ} \rightarrow (1\text{-}\mathrm{Cat})$ by a classifying topos
$\mathbb{B} = \mathbb{B}_{C_n}$ , and then to find a 2-universal $\Pi_n(\mathbb{B}) \xrightarrow{\ \phi\ } C_n$ in the sense that
for all $T$ , the functor
\[ \mathit{Hom}_{(\mathrm{Top})}(T, \mathbb{B}) \longrightarrow \tau_1 \mathit{Hom}(\Pi_n(T), C_n) \]
\[ u \longmapsto \phi_0 \Pi_n(u) \]
is an equivalence. For $n = 1$ one knows that the usual classifying topos
of $C_1$ does the job, but for $n = 2$ already, I doubt that this universal
problem has a solution. This is perhaps related to the fact that the
``van Kampen Theorem'', which one can express by saying that the 2-functor
$T \longmapsto \Pi_1(T)$ of locally 1-connected toposes to groupoids transforms (up to
1-equivalence) amalgamated sums to amalgamated sums (and more generally commutes
with inductive 2-limits), is doubtless no longer true for $\Pi_2(T)$ . Thus,
if $T$ is a topological space which is the union of two closed sets, $T_1$ and
$T_2$ , it is doubtless not true that giving a locally constant 1-champ
on $T$ ``is equivalent to'' giving a locally constant 1-champ $F_i$ on $T_i$

\mspage{124}
